% ==============================================================
% Kapitel 6: Systemarchitektur
% ==============================================================

\section{Systemarchitektur}

% -------------------------------------------------------------------
\subsection{Überblick}
\label{sec:systemarchitektur_ueberblick}

Das \gls{ZEC}-Timingsystem besteht aus zwei zentralen Teilen: einer webbasierten Anwendung,
welche zur Anzeige der Renndaten und zur Verwaltung der Teams und Fahrer dient, sowie einem
Backend-Server, welcher als Microservice-Architektur umgesetzt ist.
Der Server übernimmt die gesamte Geschäftslogik, die Authentifizierung, 
die Datenverwaltung und die Kommunikation mit der Webanwendung und anderen externen Systemen.
Abbildung~\ref{fig:systemarchitektur} zeigt einen Überblick über alle 
Komponenten des Servers und deren Zusammenspiel intern sowie nach außen.

\begin{figure}[h]
    \centering
    \includegraphics[width=14cm]{systemarchitektur.png}
    \caption{Systemarchitektur}
    \label{fig:systemarchitektur}
\end{figure}

Alle Anfragen von außen laufen über ein zentrales 
\textbf{Nginx API-Gateway}, welches als einziger Einstiegspunkt des Systems 
fungiert. Es übernimmt die \gls{JWT}-Validierung jeder eingehenden Anfrage 
über das \texttt{auth\_request}-Modul und leitet diese anschließend 
an den zuständigen Microservice weiter. So ist sichergestellt, 
dass kein Service ohne gültige Authentifizierung erreichbar ist.

Die Authentifizierung selbst wird vom \textbf{Auth Service} in 
Zusammenarbeit mit \textbf{Keycloak} übernommen. Keycloak stellt als 
dedizierter Identity Provider die \gls{JWT} aus und verwaltet 
Benutzerkonten sowie Rollen. Der Auth Service fungiert dabei als 
Vermittler zwischen Keycloak und den restlichen Services und 
synchronisiert Benutzerdaten zwischen \textbf{Keycloak} und \textbf{User Service}.

Die eigentliche Geschäftslogik ist auf fünf spezialisierte 
\textbf{Business Services} aufgeteilt:

\begin{itemize}
    \item \textbf{User Service:} Verwaltung der Benutzerkonten und 
          Teamzugehörigkeiten
    \item \textbf{Team Service:} Speicherung und Verwaltung der 
          teilnehmenden Teams und Fahrerinformationen
    \item \textbf{Challenge Service:} Verwaltung der einzelnen Bewerbe 
          und deren Konfigurationen
    \item \textbf{Attempt Service:} Erfassung der Zeitversuche und 
          Auslösen der Punkteberechnung
    \item \textbf{Score Service:} Berechnung der Gesamtpunktestände 
          und Bereitstellung der Ergebnislisten
\end{itemize}

Jeder dieser Services verfügt über eine eigene dedizierte 
\textbf{PostgreSQL-Datenbank}, wodurch eine klare Datentrennung 
gemäß dem Microservice-Prinzip gewährleistet wird und Services 
unabhängig voneinander skaliert werden können. Keycloak betreibt 
ebenfalls eine eigene Datenbankinstanz zur Verwaltung von 
seiner Daten.

Die gesamte Infrastruktur ist mittels \textbf{Docker} containerisiert 
und kann als Docker-Compose-Stack betrieben werden, was eine flexible 
und reproduzierbare Bereitstellung auf unterschiedlicher Hardware 
ermöglicht.

\newpage

% -------------------------------------------------------------------
%\subsection{Server}
% Autor: David | Umfang: 8–10 Seiten
% TODO: Alle Microservices der Reihe nach erklären
%       + Datenbankschema (max. 3 Tabellen pro Service)
%       + Web App


% -------------------------------------------------------------------
\subsection{Web App}
% Autor: Niklas | Umfang: 2–3 Seiten
% TODO: Erklärung Kommunikation mit Server und MQTT (Zeitstempel)